% Разработать метод поддержки принятия решений в рамках выполнения пошаговой инструкции, заданной в графовом формате.
% Описать основные особенности предлагаемого метода.
% Сформулировать ограничения предметной области.  
% Изложить ключевые этапы метода в виде функциональной модели и схем алгоритмов.
% Выделить структуры данных, необходимые этому методу.

% NOTE:
% В конструкторском разделе описывается разрабатываемый метод
% В случае если в бакалаврском проекте разрабатывается новый метод или алгоритм, необходимо подробно изложить их суть, привести все необходимые для их реализации математические выкладки, обосновать последовательность этапов выполнения
% При этом для каждого этапа следует выделить необходимые исходные данные и получаемые результаты
% Для описания метода или алгоритма - Схема алгоритма
% В конце описания разработанного алгоритма должны быть приведены выбранные способы тестирования и сами тесты
% Перед формированием тестовых наборов данных целесообразно указать выделенные классы эквивалентности
% (тут же могут быть приведены выкладки по теоретическим рассчетам требуемой памяти и эффективности алгоритма; эти результаты могут быть использованы для обоснования правильности выбора метода из уже имеющихся альтернативных вариантов)
% Также должно быть приведено описание структуры разрабатываемого ПО, оно включает в себя:
% - описание общей структуры - определение основных частей (компонентов) и их взаимосвязей по управлению и по данным
% - декомпозицию компонентов и построение структурных иерархий
% - проектирование компонентов
% Для графического представления такого описания (если есть необходимость), следует использовать IDEF0 с декомпозицией исходной задачи на несколько уровней

% Рек. Объем 25-30 страниц
В данном разделе будут приведены требования и ограничения разрабатываемого метода, будут описаны основные этапы разрабатываемого метода, сценарии тестирования, классы эквивалентности тестов, а также структура разрабатываемого ПО.

\section{Требования и ограничения метода}

Метод поддержки принятия решений в рамках выполнения пошаговой инструкции, заданной в графовом формате должен:
\begin{enumerate}
    \item работать с описаниями инструкций, заданных в описанном графовом формате;
    \item хранить текущий прогресс пользователя при выполнении инструкции;
    \item уметь выбирать лучший вариант действия при наличии нескольких вариантов следующих действий;
    \item выполнять переход по графу в соответствие с действиями пользователя;
    \item предоставлять пользователю рекомендации при объявлении пользователю выбранного варианта следующих действий;
    \item из возможных вариантов рекомендаций выбирать наиболее подходящие;
    \item для корректировки весов давать возможность экспертам корректировать алгоритм рекомендаций.
\end{enumerate}

% Возможно дебаг экспертом стоит оставить на позжую часть работы
% Также можно разделить алгоритмы рекомендации и движение по графу

\section{Описание разрабатываемого метода}

Поставленную задачу можно разбить на следующие этапы:
\begin{enumerate}
    \item планирование следующих действий;
    \item формирование рекомендаций для оператора;
    \item вывод следующего действия и соответствующих рекомендаций оператору;
    \item обработка результатов действия пользователя.
\end{enumerate}

Алгоритм формирования рекомендаций будет описан в данном разделе далее.

Основные этапы разрабатываемого метода представлены на IDEF0 диаграмме первого уровня (см. Рисунок ).

Рисунок

\section{Планирование следующих действий}

Исходные данные --- инструкция в графовом формате и вершина графа, описывающая текущий этап выполнения инструкции.
Получаемый результат --- предлагаемая следующая вершина графа и переход, описывающий действие, необходимое для достижения нужного состояния.

На рисунке представлена IDEF0 диаграмма планирования следующих действий.

% Выбор кратчайшего пути с помощью чего-то вроде волнового алгоритма 
% или
% Предварительная разметка экспертами главного пути
% Можно попытаться скрыть выбор следующего действия под словом планирование


% Возможно стоит поменять формат графа
Текущее положение описывается следующими сущностями:
\begin{enumerate}
    \item possible\_actions --- массив, $i$-й элемент которого описывает $i$-ый из возможных из текущего состояния переходов в $i$-ое состояние;
    \item linked\_states --- массив, $i$-й элемент которого описывает $i$-ый все состояния, в которые возможны переходы;
    \item actor --- субъект, выполняющий действие;
    \item object --- объект, над которым выполняется действие.
\end{enumerate}

% Если мы может использовать один переход для нескольких состояний/сущностей, где-то будет необходимо хранить актуальные начало-конец, а не только возможные. 
% Если для каждого объекта они свои, то скорее всего придется тут. 
% Если они общие (все начинают ненарезанные -> все становятся нарезанными), то такое необходимости нет, объекты хранятся на уровне выше.
Каждое переход описывается следующими данными:
\begin{enumerate}
    \item text --- текст, описывающий действие пользователя;
    \item tags --- массив, $i$-й элемент которого описывает тэг, описывающий данное действие;
    \item starting\_state --- состояние, в котором изначально может находиться объект действия;
    \item ending\_state --- состояние, в котором конечно может находится объект.
\end{enumerate}

Tag описывается следующими данными:
\begin{enumerate}
    \item name --- текстовое поле, описывает название;
    \item coefficients --- массив, $i$-й элемент которого является кортежем типа $(сriteria\_name, weight, C)$, где criteria\_name --- имя критерия, weight --- значение типа float от 0 до 1, описывающее важность критерия для данного тэга;
    \item recommendation --- текстовое поле, описывающее рекомендацию, которую необходимо дать оператору при присутствии tag у действияю.
\end{enumerate}



Каждый 

% Редактор графа(?) + визуализатор возможно(граф виз) - ?
% Подгрузка общих кусков графа на случай возможных общих ситуаций - +
%
 